\subsection{Injective, Surjective, and Bijective Maps}

\begin{definition} \label{def:injective-surjective-bijective} \hlt{Injective, Surjective, and Bijective Maps}\\
Let $X,Y$ be sets and let $f:X\to Y$.
\begin{enumerate}[label=(\roman*)]
\setlength{\itemsep}{0pt}
\item \hlt{Injective}: $(\forall x_1,x_2\in X)\,[f(x_1)=f(x_2)\Longrightarrow x_1=x_2]$.
\item \hlt{Surjective}: $\operatorname{range}(f)=Y$.
\item \hlt{Bijective}: $f$ is both injective and surjective.
\end{enumerate}
\end{definition}

\begin{illustrationfigure}[H]
  \centering
  \begin{tikzpicture}[
  scale=0.80,
  >=Stealth,
  every node/.style={font=\scriptsize},
  set/.style={draw=gray!65, fill=gray!10, line width=0.45pt},
  point/.style={circle, fill=black, inner sep=1.05pt},
  missed/.style={circle, fill=white, draw=black, line width=0.45pt, inner sep=1.05pt},
  map/.style={->, gray!70, line width=0.48pt}
]
  \foreach \xshift in {0,4.2,8.4} {
    \begin{scope}[xshift=\xshift cm]
      \draw[set] (0,0) ellipse (0.6 and 1.15);
      \draw[set] (2.9,0) ellipse (0.6 and 1.15);
      \node at (0,-1.42) {$X$};
      \node at (2.9,-1.42) {$Y$};
    \end{scope}
  }

  \begin{scope}
    \coordinate (a1) at (0,0.72);
    \coordinate (a2) at (0,0.18);
    \coordinate (a3) at (0,-0.36);
    \coordinate (a4) at (0,-0.86);
    \coordinate (b1) at (2.9,0.58);
    \coordinate (b2) at (2.9,0.00);
    \coordinate (b3) at (2.9,-0.62);
    \foreach \p in {a1,a2,a3,a4,b1,b2,b3} \node[point] at (\p) {};
    \draw[map] (a1) -- (b1);
    \draw[map] (a2) -- (b2);
    \draw[map] (a3) -- (b2);
    \draw[map] (a4) -- (b3);
  \end{scope}

  \begin{scope}[xshift=4.2cm]
    \coordinate (a1) at (0,0.68);
    \coordinate (a2) at (0,0.02);
    \coordinate (a3) at (0,-0.68);
    \coordinate (b1) at (2.9,0.78);
    \coordinate (b2) at (2.9,0.24);
    \coordinate (b3) at (2.9,-0.28);
    \coordinate (b4) at (2.9,-0.82);
    \foreach \p in {a1,a2,a3,b1,b2,b3} \node[point] at (\p) {};
    \node[missed] at (b4) {};
    \draw[map] (a1) -- (b1);
    \draw[map] (a2) -- (b2);
    \draw[map] (a3) -- (b3);
  \end{scope}

  \begin{scope}[xshift=8.4cm]
    \coordinate (a1) at (0,0.62);
    \coordinate (a2) at (0,0.00);
    \coordinate (a3) at (0,-0.62);
    \coordinate (b1) at (2.9,0.62);
    \coordinate (b2) at (2.9,0.00);
    \coordinate (b3) at (2.9,-0.62);
    \foreach \p in {a1,a2,a3,b1,b2,b3} \node[point] at (\p) {};
    \draw[map] (a1) -- (b1);
    \draw[map] (a2) -- (b2);
    \draw[map] (a3) -- (b3);
  \end{scope}
\end{tikzpicture}

  \caption{Left: surjective with a collision; center: injective but not surjective; right: bijective.}
  \label{fig:injective-surjective-bijective-maps}
\end{illustrationfigure}

Figure \ref{fig:injective-surjective-bijective-maps} contrasts the three conditions in Definition \ref{def:injective-surjective-bijective}. The left map is surjective but not injective, whereas the center map is injective but not surjective. The right map is bijective.

\begin{example} \label{ex:square-map-image-preimage} \hlt{Image and Preimage for the Square Map}\\
Define $f:\R\to\R$ by $f(x)\defeq x^2$. If $x\in[-2,1]$, then $0\leq f(x)\leq4$. Conversely, if $y\in[0,4]$, then $-\sqrt{y}\in[-2,1]$ and $f(-\sqrt{y})=y$.\\
Hence $f([-2,1])=[0,4]$. For $x\in\R$, $x\in f^{-1}([1,4])\Longleftrightarrow1\leq x^2\leq4\Longleftrightarrow1\leq\abs{x}\leq2\Longleftrightarrow x\in[-2,-1]\cup[1,2]$.\\
Since $f(-1)=f(1)$ and $-1\neq1$, the map is not injective. Since $(\forall x\in\R)\,f(x)\geq0$, one has $-1\notin\operatorname{range}(f)$, so the map is not surjective.\\
\hlt{Key Insight}: Repeated values obstruct injectivity; omitted codomain values obstruct surjectivity.
\end{example}
